\subsection{半单环上的模}

命题 4. —— 设 A 是一个环。下列性质等价：

\begin{enumerate}
\def\labelenumi{(\roman{enumi})}
\item
  环 A 是半单的。
\item
  每个 A-模都是半单的。
\item
  存在一个生成且半单的 A-模。
\item
  存在一个忠实半单且反模有限生成的 A-模。
\item
  每个 A-模都是投射的。
\item
  每个单生成的 A-模都是投射的。
\end{enumerate}

*关于半单环的其他刻画，参见 X, §8, n° 4, 第 140 页的命题 6。*

先证明 (i) 蕴含 (ii) 与 (v)。设环 A 是半单的，并考虑一个左 A-模 M。按假设，A-模 \(A_{s}\) 是半单的，故每个自由 A-模都是半单的。由 II, §1, No.~11, 第 218 页的命题 20，存在一个自由 A-模 L 和一个从 L 到 M 的满 A-线性映射 u。设 N 是 u 的核。由于 A-模 L 是半单的，存在 L 中与 N 互补的半单子模 \(N'\)（VIII, 第 56 页，定理 1）。A-模 \(N'\) 是投射的，且 u 诱导出从 \(N'\) 到 M 的同构。因此，M 是半单的且投射的。

\begin{enumerate}
\def\labelenumi{(\roman{enumi})}
\setcounter{enumi}{1}
\item
  \(\Rightarrow\) (iii)：若每个 A-模都是半单的，则 A-模 \(A_{s}\) 是半单的；它还是生成模。
\item
  \(\Rightarrow\) (iv)：若 M 是一个生成 A-模，则它是忠实的（VIII, 第 80 页，定理 1 的推论），且它的反模是有限生成的（VIII, 第 99 页，推论 1）。
\item
  \(\Rightarrow\) (i)：设 M 是一个忠实半单且反模有限生成的 A-模。由 VIII, 第 8 页的引理 4，存在一个自然数 m，使得 \(A_{s}\) 同构于 \(M^{m}\) 的一个子模。A-模 \(M^{m}\) 是半单的，因而 \(A_{s}\) 亦然。
\end{enumerate}

蕴含关系 \((v) \Rightarrow (vi)\) 是显然的。

\begin{enumerate}
\def\labelenumi{(\roman{enumi})}
\setcounter{enumi}{5}
\tightlist
\item
  \(\Rightarrow\) (i)：设每个单生成的 A-模都是投射的。设 \(\mathfrak{a}\) 是 A 的一个左理想。由于 A-模 \(A_s / \mathfrak{a}\) 是投射的，存在 A 的一个左理想 \(\mathfrak{b}\)，使得 \(A_s\) 是 \(\mathfrak{a}\) 与 \(\mathfrak{b}\) 的直和（II, §2, No.~2, 第 231 页，命题 4）。于是环 A 是半单的（VIII, 第 135 页，定理 1）。
\end{enumerate}

引理 1. —— 设 A 是一个左阿廷环，M 是一个单 A-模。则环 \(A_{M}\) 是单的。

由于环 A 是左阿廷的，由 VIII, 第 7 页的命题 5，环 \(A_{M}\) 亦然。而 M 是一个忠实单 \(A_{M}\)-模，故环 \(A_{M}\) 是单的（VIII, 第 119 页，命题 1）。

命题 5. —— 设环 A 是半单的。设 M 是一个左 A-模。M 的反模是有限生成的，且我们有 \(A_{M} = A_{M}^{\prime\prime}\)。

先考虑 M 是单 A-模的情形。由引理 1，环 \(A_{M}\) 是单的。于是命题 5 由 VIII, 第 120 页的引理 1 得出。

现在考虑一般情形。单 A-模的类所成的集 S 是有限的（VIII, 第 136 页，命题 1）。对每个 \(\lambda \in S\)，选取一个 A-模 \(S_{\lambda}\)（类为 \(\lambda\)），并把 \(S_{\lambda}\) 的交换子的反体记作 \(D_{\lambda}\)。由引理 1 应用于单环 \(A_{S_{\lambda}}\)，\(S_{\lambda}\) 是体 \(D_{\lambda}\) 上的一个有限维向量空间；用 \(m(\lambda)\) 表示它的维数。设 B 是 M 的自同态环的反环。我们在 VIII, 第 84 页已看到，存在作为 B-模都是单的 \((D_{\lambda}, B)\)-双模 \(V_{\lambda}\)，以及一个 \((A, B)\)-双模同构（从 M 到 \(\oplus_{\lambda \in S} S_{\lambda} \otimes_{D_{\lambda}} V_{\lambda}\)）。作为 B-模，M 同构于 \(\oplus_{\lambda \in S} V_{\lambda}^{m(\lambda)}\)。由于 S 与诸 \(m(\lambda)\) 都是有限的，M 是一个有限生成的 B-模。

A-模 M 是半单的，且它的反模是有限生成的。于是由 VIII, 第 83 页的命题 4，我们有 \(A_{M} = A_{M}^{\prime\prime}\)。

命题 6. —— 设 M 是一个有限生成的半单 A-模。用 \(\mathcal{S}_{\mathrm{M}}\) 表示它的支集（VIII, 第 66 页），用 B 表示它的自同态环。对每个 \(\lambda \in \mathcal{S}_{\mathrm{M}}\)，选取一个单 A-模 \(S_{\lambda}\)（类为 \(\lambda\)），并把左 B-模 \(\mathrm{Hom}_{\mathrm{A}}(S_{\lambda}, \mathrm{M})\) 记作 \(V_{\lambda}\)。

\begin{enumerate}
\def\labelenumi{\alph{enumi})}
\item
  环 B 是半单的。
\item
  映射 \(\lambda \mapsto \operatorname{cl}(V_{\lambda})\) 是从 M 的支集 \(\mathcal{S}_{\mathrm{M}}\) 到单 B-模的类所成的集上的双射。
\item
  对每个 \(\lambda \in \mathcal{S}_{\mathrm{M}}\)，A-模 M 的 \(\lambda\) 型同型分量等于 B-模 M 的 \(V_{\lambda}\) 型同型分量。
\end{enumerate}

作为 B-模，M 是半单的（VIII, 第 85 页，命题 5）且忠实的。由于我们有 \(A_{M} \subset \operatorname{End}_{B}(M)\)，它的反模是有限生成的。于是环 B 是半单的（命题 4）。若 \((x_{1}, \ldots, x_{r})\) 是 A-模 M 的一个生成序列，则映射 \(b \mapsto (bx_{1}, \ldots, bx_{r})\)（从 \(B_{s}\) 到 \(M^{r}\)）是 B-线性的且单的。每个单 B-模都同构于 \(B_{s}\) 的一个子模（VIII, 第 136 页，命题 1），从而同构于 M 的一个 B-子模。于是本命题随即由 VIII, 第 85 页的命题 5 得出。

命题 7. —— 设 A 是一个半单环。

\begin{enumerate}
\def\labelenumi{\alph{enumi})}
\item
  每个有限生成的 A-模都是自反的（II, §2, No.~7, 第 239 页）。
\item
  对每个单左 A-模 S，S 的对偶右 A-模 \(S^{*}\) 是单的，且映射 \(\lambda \mapsto \operatorname{cl}(\lambda^{*})\) 定义了从单 A-模的类所成的集到单右 A-模的类所成的集上的一个双射。
\item
  设 M 是一个有限生成的左 A-模。M 的对偶右 A-模 \(M^{*}\) 是有限生成的且与 M 具有相同的长度。此外，对每个单 A-模 S，我们有 \([M:S]=[M^{*}:S^{*}]\)。
\end{enumerate}

设 \(\mathrm{M}\) 是一个有限生成的 A-模。

由 VIII, 第 138 页的命题 4，A-模 M 是有限生成的且投射的；因而由 II, §2, No.~7, 第 240 页的推论 4，它是自反的。特别地，每个单 A-模都是自反的。由此还可得出，两个有限生成的左模同构当且仅当它们的对偶同构。

设 S 是一个单左 A-模。设 \((\mathrm{T}_{i})_{i\in\mathrm{I}}\) 是一个以 S 的对偶 \(S^{*}\) 为直和的单右 A-模族。由于 S 是自反的，它同构于 \((\mathrm{S}^{*})^{*}\)，从而同构于 \(\prod_{i\in I}T_{i}^{*}\)。诸模 \(T_{i}\) 中的每一个都是自反的；特别地，我们有 \(T_{i}^{*}\neq0\)（对每个 \(i\in I\)）。由于单模 S 同构于 \(\prod_{i\in I}T_{i}^{*}\)，集 I 只有一个元素，故 \(S^{*}\) 是单的。

由于 M 是半单的且有限生成的，它是单子模 \(S_{1},\ldots,S_{r}\) 的直和。于是 \(M^{*}\) 同构于族 \(S_{1}^{*},\ldots,S_{r}^{*}\) 的直和，而我们刚看到诸模 \(S_{i}^{*}\) 都是单的。由此立即得到断言 c)。

